\chapter{Discussion}
\label{c:discussion}

% Struktur
\section{Threats to Validity}

\subsection{Insufficient Reflection of real-world Vulnerabilities}
A point of criticism often raised when discussing intentionally vulnerable Android apps, such as the ones found in the MASTG reference apps catalog, is about the extent to which these apps accurately reflect the vulnerabilities found in real-world Android apps that can be downloaded from the Google Playstore and other platforms. Especially when using them for benchmarking SAST and DAST tools, which are then used for testing real-world apps, it is important for the benchmarking apps to closely mimic real-world ones. If this is not the case, they lose their function for evaluating such security tools, since a low or high score on the reference apps would not necessarily translate to the tools performances on real apps. The reference apps might include vulnerabilities which are rarely seen in real apps, while missing out on crucial vulnerabilities that a tool has to be able to detect. 

Speaking against this point of criticism is the relevancy and status of the OWASP MASTG as an industry standard for evaluating mobile application security. There are numerous papers \cite{paper_01, paper_02, paper_04, paper_05, paper_07, paper_14} that show the effectiveness of the MASTG when it comes to finding vulnerabilities. \cite{paper_03} even found that using the MASTG leads to finding more vulnerabilities in real world apps than using the automated testing tools MobSF. \todo{add source for MobSF??} Since the MASTG deals with the same vulnerabilities as discussed in the MASWE and implemented in the MASTG reference apps, this backs the legitimacy and real world application the vulnerabilities implemented and found in MASTG reference apps have. For this to be case however, it is crucial that the MASTG reference apps are under constant maintenance and evolving alongside the OWASP MAS project, which in turn adapts to the security issues most prevalent at the time. \todo{word better than "at the time"}

\subsection{Insufficient Complexity of Vulnerabilities}
Besides not reflecting the vulnerabilities found in real world apps, it can also be that the vulnerabilities implemented in intentionally vulnerable apps such as the MASTG reference apps are of the right type, but are implemented in a shallow and unrealistic way, which fails to reflect how the vulnerabilities look like in real world apps. This would lead to mobile app security testing tools performing well on the MASTG reference apps due to the easy to spot nature of their implemented vulnerabilities, but struggle to do the same on real world apps that include similar yet more complex vulnerabilities. 

To counter this argument a study would have to be conducted, collecting vulnerabilities from real world apps and comparing them to the ones found in the MASTG reference apps catalog. Zhu et al. \cite{paper_19} evaluated the effectiveness of 11 different SAST tools on four different benchmarks, of which two are synthetic and the other two are real-world benchmarks based on vulnerabilities found in real-world apps. One of the synthetic benchmarks is based on the OWASP MASTG reference apps catalog. The authors find that there is no significant difference in the performance of the 11 SAST tools between the synthetic and the real-world benchmarks. This illustrates that the vulnerabilities implemented in MASTG reference apps are similar both in types as well as in complexity and depth to the ones found in real-world apps.

\subsection{Relevance of even MASWE Vulnerabilities Coverage}
While selecting which MASWE vulnerabilities to implement, we focused on vulnerabilities which have not been implemented by other MASTG reference apps. Our rationale behind this decision is that all vulnerabilities as described by OWASP MASWE are relevant and should have some form of representation in the reference apps catalog, both for educational purposes and for evaluating security testing tools. An argument can be made against the validity of this approach, suggesting to focus on the more often implemented vulnerabilities instead of the opposite group. It can be argued that the vulnerabilities that have been implemented by multiple MASTG existing reference apps are the ones. \todo{SInce they are more popular because they are more important and more often found yada yada is the argument}

\section{Open Problems}
